\chapter{Metodologia}

Este capítulo descreve o caminho metodológico adotado para transformar o tema do
avaliador de serviços públicos em uma proposta de solução. A metodologia foi
organizada em três partes: metodologia do trabalho, metodologia de
desenvolvimento da solução e solução de análise e desenvolvimento de dados.

\section{Metodologia do Trabalho}

Este TCC possui natureza aplicada, pois busca construir uma solução computacional
para um problema relacionado à melhoria de serviços públicos digitais. Quanto aos
objetivos, a pesquisa é exploratória, uma vez que parte de modelos e normas já
existentes para propor um artefato inicial, ainda sujeito a validação e
refinamento.

A condução do trabalho segue uma abordagem inspirada em Design Science Research.
Segundo \citeonline{hevner2004designscience}, pesquisas em sistemas de
informação podem produzir conhecimento por meio da construção e avaliação de
artefatos, como modelos, métodos, arquiteturas e sistemas. A proposta deste TCC
se enquadra nessa lógica porque o resultado principal é um protótipo de API e
interface que materializa uma forma de coletar e analisar sinais de dificuldade
em serviços públicos digitais.

As etapas metodológicas adotadas foram:

\begin{enumerate}
    \item revisão de bases legais, modelos de avaliação de serviços públicos,
    conceitos de qualidade percebida, satisfação, usabilidade e governo digital;
    \item definição do problema e dos objetivos com foco em coleta de feedback de
    baixo esforço e identificação de sinais de dificuldade durante a navegação;
    \item elaboração do fluxo conceitual da solução, contemplando registro de
    cliques, análise de tempo, sugestão de redirecionamento e coleta de feedback;
    \item construção de um protótipo funcional com API, páginas simuladas de
    serviço público, formulário de avaliação e painel administrativo;
    \item organização dos dados coletados em indicadores capazes de apoiar a
    análise de sessões críticas, avaliações, NPS e características recorrentes;
    \item verificação preliminar por simulação, observando se o fluxo proposto é
    coerente com os objetivos e se o protótipo permite demonstrar a ideia central.
\end{enumerate}

A revisão teórica não teve como finalidade esgotar a literatura sobre qualidade
de serviços. Ela foi direcionada à construção da solução. Foram priorizadas
fontes normativas brasileiras, modelos de avaliação de serviços públicos, bases
de usabilidade, design centrado no usuário, métricas de serviços digitais e
referências de desenvolvimento de software.

\section{Metodologia do Desenvolvimento da Solução}

O desenvolvimento da solução foi conduzido por prototipação incremental. Essa
estratégia é adequada quando os requisitos ainda estão em amadurecimento e
quando o artefato precisa ser demonstrado rapidamente para discussão e ajuste.
Em engenharia de software, protótipos permitem explorar requisitos, reduzir
incertezas e comunicar alternativas de solução antes de uma implementação
definitiva \cite{Sommerville2007}.

O protótipo foi organizado em quatro componentes principais:

\begin{itemize}
    \item \textbf{páginas públicas simuladas}: telas que representam uma página
    inicial de órgão público e páginas de serviço, com estrutura visual mais
    próxima de portais governamentais reais;
    \item \textbf{script cliente}: código executado no navegador para iniciar uma
    sessão, registrar cliques, registrar tempo de permanência e enviar eventos
    para a API;
    \item \textbf{API de avaliação}: backend responsável por receber eventos,
    atualizar métricas da sessão, aplicar regras de dificuldade e devolver ações
    para o frontend;
    \item \textbf{interfaces de retorno}: popup de sugestão de redirecionamento,
    página de feedback e painel administrativo com dados agregados.
\end{itemize}

O fluxo funcional definido para o protótipo parte do acesso do usuário à página
inicial do órgão ou à página de um serviço. A partir desse acesso, o script
cliente inicia uma sessão anônima ou identificada, registra eventos de clique e
envia os dados para a API. A API atualiza métricas da sessão, como quantidade de
cliques, páginas visitadas e tempo de permanência. Quando cliques ou tempo ficam
acima da faixa esperada, a dificuldade é registrada silenciosamente para o
administrador.

Serviços marcados como prioritários ou de alto tráfego recebem tratamento
adicional. Se o usuário apresenta sinais de dificuldade procurando um serviço
desse tipo, a API pode retornar uma ação de popup sugerindo o redirecionamento
para a página correta. Caso o usuário aceite, o frontend realiza o
redirecionamento; caso recuse, a navegação continua. Para evitar interrupções
repetitivas, o popup deve ser exibido apenas uma vez por sessão, serviço e motivo
de acionamento.

Quando a sessão apresenta dificuldade ou tempo acima da média, a API pode
retornar uma ação para abrir a página de feedback. O formulário utiliza uma
escala de estrelas para avaliação geral, seleção de características do problema
e uma pergunta de NPS. A decisão por estrelas e características objetivas busca
reduzir o esforço de resposta e, ao mesmo tempo, qualificar a causa da
insatisfação. A pergunta de NPS foi incluída como indicador complementar de
recomendação, tomando como base a simplicidade do método proposto por
\citeonline{reichheld2003nps}.

\section{Solução de Análise e Desenvolvimento de Dados}

A solução de dados foi definida para responder a duas perguntas administrativas:
onde o usuário demonstra dificuldade e quais aspectos do serviço precisam ser
priorizados. Para isso, o protótipo combina dados comportamentais, coletados
durante a navegação, e dados declarativos, informados pelo usuário no feedback.

Os dados comportamentais incluem identificador de sessão, página acessada,
serviço associado, origem do evento, instante, coordenadas aproximadas do clique,
quantidade acumulada de cliques e tempo de permanência. Esses dados permitem
detectar sessões com esforço acima do esperado. Os dados declarativos incluem
nota por estrelas, características selecionadas, NPS e comentário opcional. Essa
combinação segue a ideia de que satisfação e qualidade percebida devem ser
interpretadas por atributos, e não apenas por uma nota geral
\cite{silva2024atributosApi,iso9241112018}.

\begin{table}[h]
\centering
\caption{Dados previstos no protótipo}
\label{tab:dados-prototipo}
\begin{tabular}{p{3.4cm} p{5.0cm} p{5.0cm}}
\hline
\textbf{Grupo de dados} & \textbf{Exemplos} & \textbf{Uso esperado} \\
\hline
Navegação & Sessão, página, serviço, clique, tempo e origem & Detectar esforço
acima da faixa esperada e registrar dificuldade silenciosa \\
Feedback & Estrelas, atributos, NPS e comentário opcional & Explicar a percepção
do usuário e identificar pontos recorrentes de melhoria \\
Administração & Alertas, sessões críticas, médias, distribuições e recorrência
de atributos & Apoiar priorização, acompanhamento e tomada de decisão \\
\hline
\end{tabular}
\end{table}

As regras de análise do protótipo são configuráveis. Cada serviço pode possuir
uma faixa esperada de cliques e tempo, além de uma marcação de prioridade ou alto
tráfego. Quando a faixa é ultrapassada, o sistema registra um alerta. Em serviços
comuns, a navegação continua sem popup. Em serviços prioritários, a API pode
sugerir redirecionamento. Essa diferenciação evita que todo sinal de dificuldade
gere interrupção para o usuário.

O painel administrativo deve apresentar os resultados de forma agregada. Entre os
indicadores previstos estão: total de sessões, sessões com dificuldade, serviços
com maior quantidade de alertas, média de estrelas, distribuição das notas, NPS,
atributos mais selecionados e comentários recentes. A lógica acompanha a
orientação de que métricas de serviços digitais devem ser usadas para acompanhar
desempenho e identificar oportunidades de melhoria \cite{govuk2021usingData}.

Do ponto de vista de privacidade, a solução foi pensada para trabalhar com
sessões técnicas e dados agregados. A coleta de eventos deve ser limitada ao que
é necessário para analisar a experiência do serviço, respeitando princípios de
finalidade, adequação e necessidade previstos na LGPD \cite{brasil2018lgpd}. Em
uma implantação real, seriam necessários mecanismos adicionais, como aviso de
coleta, governança de dados, política de retenção, controle de acesso e avaliação
jurídica.

\section{Base de Fundamentação da Solução}

A resposta à pergunta ``com o que a solução se baseou?'' pode ser organizada em
cinco grupos de referência. O primeiro grupo é normativo: Lei
n\textsuperscript{o} 13.460/2017, Decreto n\textsuperscript{o} 9.094/2017, Lei
n\textsuperscript{o} 14.129/2021 e LGPD, que justificam avaliação, participação,
governo digital e cuidado com dados. O segundo grupo é técnico-governamental:
BRASP, Qualidade Gov.br, modelo unificado e estudo dos atributos da API de
avaliação, que orientam qualidade percebida, satisfação, efetividade e atributos
de serviço público.

O terceiro grupo é de experiência do usuário: ISO 9241-11 e ISO 9241-210, usadas
para fundamentar usabilidade e design centrado no usuário. O quarto grupo é de
avaliação de sistemas e serviços digitais: DeLone e McLean e o Service Manual do
GOV.UK, que sustentam o uso de métricas de uso, satisfação, conclusão e melhoria
contínua. O quinto grupo é de desenvolvimento de software e pesquisa aplicada:
Sommerville e Design Science Research, usados para justificar a construção de um
protótipo como artefato de pesquisa.

Assim, a proposta não se baseia apenas na comparação com aplicativos privados.
As avaliações por estrelas, motivos e NPS são referências de interação de baixo
esforço; a finalidade, os atributos, a análise dos dados e os cuidados de
governança são fundamentados em avaliação de serviços públicos, governo digital,
usabilidade, engenharia de software e proteção de dados.
